\chapter{Decorators, Context Managers, and the 2.x Twilight (Python 2.4--2.7)}

The late 2.x line added three features that turned Python's first-class functions and frames
into reusable structure: \textbf{decorators} (callable transformers applied with \texttt{@}),
\textbf{context managers} (the \texttt{with} protocol for deterministic acquire/release), and
\textbf{generator coroutines} (two-way \texttt{send}/\texttt{throw}/\texttt{close}). Decorators
are function composition made declarative; context managers are Python's answer to C++ RAII
built on \texttt{try}/\texttt{finally}; and generator coroutines are the suspendable frame
that, a decade later, became \texttt{async}/\texttt{await}. This chapter teaches all three at
working depth and grounds each in the bytecode the 3.13 interpreter actually emits.

\section*{Section Index}
\begin{itemize}
  \item \textbf{5.1} Decorators: composition made declarative (PEP 318, PEP 3129)
  \item \textbf{5.2} Context managers and \texttt{with} (PEP 343)
  \item \textbf{5.3} Generator coroutines: \texttt{send}, \texttt{throw}, \texttt{close} (PEP 342)
  \item \textbf{5.4} Performance and anti-patterns
  \item \textbf{5.5} Summary and cross-references
\end{itemize}

\section{Decorators: composition made declarative (PEP 318, PEP 3129)}

\textbf{Why this exists.} Before Python 2.4, wrapping a function meant defining it and then
rebinding the name, with the transformation written after and away from the definition
(\texttt{query = transaction(log(query))}). \textbf{PEP 318} introduced \texttt{@decorator}
syntax so the transformation sits \emph{on} the definition. It is pure compile-time sugar:
\texttt{@d} above \texttt{def f} means ``after building \texttt{f}, replace it with
\texttt{d(f)}.'' Stacked decorators apply \textbf{bottom-up} (nearest the \texttt{def} first).

\begin{lstlisting}[language=Python, caption={Stacked decorators compose bottom-up; functools.wraps preserves identity.}]
import functools

def dec1(f):
    @functools.wraps(f)
    def wrapper(*args, **kwargs):
        return "dec1(" + f(*args, **kwargs) + ")"
    return wrapper

def dec2(f):
    @functools.wraps(f)
    def wrapper(*args, **kwargs):
        return "dec2(" + f(*args, **kwargs) + ")"
    return wrapper

@dec1
@dec2
def target():
    return "target"

print("result:", target())
print("__name__ preserved:", target.__name__)
\end{lstlisting}

Verified output (CPython 3.13.5):

\begin{lstlisting}[caption={Output}]
result: dec1(dec2(target))
__name__ preserved: target
\end{lstlisting}

The result is the bottom-up composition $\text{target} = \text{dec1}(\text{dec2}(\text{target}))$.
The bytecode makes the mechanism explicit:

\begin{lstlisting}[caption={dis of stacked decorators on CPython 3.13.5}]
LOAD_NAME    0 (dec1)
LOAD_NAME    1 (dec2)
LOAD_CONST   0 (<code object t>)
MAKE_FUNCTION
CALL         0          # dec2(t)
CALL         0          # dec1(dec2(t))
STORE_NAME   2 (t)
RETURN_CONST 1 (None)
\end{lstlisting}

This is current bytecode: \texttt{MAKE\_FUNCTION} now takes no inline argument, and the call
opcode is the unified \texttt{CALL}, not the pre-3.11 \texttt{CALL\_FUNCTION}.

\textbf{\texttt{functools.wraps} is not optional.} A naive wrapper replaces the function's
identity---\texttt{\_\_name\_\_}, \texttt{\_\_doc\_\_}, \texttt{\_\_qualname\_\_},
\texttt{\_\_wrapped\_\_}, \texttt{\_\_module\_\_}. \texttt{@functools.wraps(f)} copies them so
introspection, tracebacks, and tools still see the original. Omitting it is the most common
decorator bug.

\textbf{Decorators with arguments} are decorator \emph{factories}---a function returning a
decorator:

\begin{lstlisting}[language=Python, caption={A parameterized decorator is a function that returns a decorator.}]
import functools

def repeat(n):
    def decorator(f):
        @functools.wraps(f)
        def wrapper(*args, **kwargs):
            return [f(*args, **kwargs) for _ in range(n)]
        return wrapper
    return decorator

@repeat(3)
def greet(name):
    return f"hi {name}"

print(greet("ada"))
\end{lstlisting}

Verified output (CPython 3.13.5):

\begin{lstlisting}[caption={Output}]
['hi ada', 'hi ada', 'hi ada']
\end{lstlisting}

\textbf{Class decorators (PEP 3129, Python 2.6)} apply the same rule to a class:
\texttt{Model = class\_decorator(Model)}. This is how \texttt{@dataclass},
\texttt{@functools.total\_ordering}, and registration decorators work (Vol IV).

\section{Context managers and \texttt{with} (PEP 343)}

\textbf{Why this exists.} Cleanup that must happen regardless of how a block exits---close a
file, release a lock, roll back a transaction---is what C++ expresses with RAII and
destructors. Python's reference-counting finalization (Chapter 2) is too imprecise to rely on
(and useless across cycles), so \textbf{PEP 343} added \texttt{with} and the
\textbf{context-manager protocol}: \texttt{\_\_enter\_\_(self)} acquires (its return value
binds to the \texttt{as} target); \texttt{\_\_exit\_\_(self, exc\_type, exc\_val, exc\_tb)}
releases. On a clean exit all three are \texttt{None}; on an exception they carry its details,
and \textbf{returning a truthy value suppresses the exception}.

\begin{lstlisting}[language=Python, caption={\_\_exit\_\_ sees the exception and can suppress it by returning truthy.}]
class Ctx:
    def __enter__(self):
        print("  __enter__")
        return self
    def __exit__(self, exc_type, exc_val, exc_tb):
        print("  __exit__ exc_type =", exc_type.__name__ if exc_type else None)
        return exc_type is ValueError      # suppress only ValueError

with Ctx() as c:
    print("  body (clean)")
print("clean exit done")

with Ctx():
    print("  body raising ValueError")
    raise ValueError("boom")
print("ValueError was suppressed by __exit__")
\end{lstlisting}

Verified output (CPython 3.13.5):

\begin{lstlisting}[caption={Output}]
  __enter__
  body (clean)
  __exit__ exc_type = None
clean exit done
  __enter__
  body raising ValueError
  __exit__ exc_type = ValueError
ValueError was suppressed by __exit__
\end{lstlisting}

\textbf{What the compiler does.} \texttt{with} desugars to a guaranteed \texttt{try}/finally
with the exception path routed through \texttt{\_\_exit\_\_}:

\begin{lstlisting}[language=Python, caption={Conceptual desugaring of: with expr as target: suite}]
mgr = expr
value = type(mgr).__enter__(mgr)
try:
    target = value          # if an "as" clause is present
    suite
except:                     # exception path
    if not type(mgr).__exit__(mgr, *sys.exc_info()):
        raise
else:                       # clean path
    type(mgr).__exit__(mgr, None, None, None)
\end{lstlisting}

The VM implements this with dedicated opcodes (\texttt{BEFORE\_WITH}/\texttt{WITH\_EXCEPT\_START})
and the frame's \textbf{exception table} (3.11+) rather than the older block stack, but the
guarantee is identical.

\begin{lstlisting}[language=Python, caption={A generator-based context manager.}]
import contextlib

@contextlib.contextmanager
def managed():
    print("  acquire")
    try:
        yield 42
    finally:
        print("  release")

with managed() as v:
    print("  using", v)
\end{lstlisting}

Verified output (CPython 3.13.5):

\begin{lstlisting}[caption={Output}]
  acquire
  using 42
  release
\end{lstlisting}

\texttt{contextlib} adds \texttt{suppress}, \texttt{closing}, \texttt{ExitStack}, and
\texttt{nullcontext}; the full reference is in Vol X. The point here is the \emph{protocol} and
that it is Python's deterministic-cleanup primitive---use it for every OS resource instead of
trusting \texttt{\_\_del\_\_}.

\section{Generator coroutines: \texttt{send}, \texttt{throw}, \texttt{close} (PEP 342)}

\textbf{Why this exists.} Python 2.5's \textbf{PEP 342} upgraded generators from one-way
producers into two-way \textbf{coroutines}---a suspended frame you can resume \emph{with a
value}, inject an exception into, or shut down. This is the conceptual seed of
\texttt{async}/\texttt{await} (Vol III). The protocol adds \texttt{.send(value)} (resume; the
paused \texttt{yield} expression evaluates to \texttt{value}; \texttt{.send(None)} equals
\texttt{next()}), \texttt{.throw(exc)} (raise at the \texttt{yield} point), and
\texttt{.close()} (raise \texttt{GeneratorExit}).

\begin{lstlisting}[language=Python, caption={A running-average coroutine driven by send(), shut down by close().}]
def averager():
    total, count, avg = 0.0, 0, None
    while True:
        try:
            x = yield avg                 # value arrives via .send(x)
        except GeneratorExit:
            print("  averager closing")
            return
        total += x
        count += 1
        avg = total / count

g = averager()
next(g)                                   # prime to the first yield
print("avg after 10:       ", g.send(10))
print("avg after 10,20:    ", g.send(20))
print("avg after 10,20,30: ", g.send(30))
g.close()
\end{lstlisting}

Verified output (CPython 3.13.5):

\begin{lstlisting}[caption={Output}]
avg after 10:        10.0
avg after 10,20:     15.0
avg after 10,20,30:  20.0
  averager closing
\end{lstlisting}

\textbf{What the interpreter actually does.} At a \texttt{yield},
\texttt{\_PyEval\_EvalFrameDefault} does not destroy the frame: it records the resume point
(\texttt{f\_lasti}), freezes the evaluation-stack depth, marks the generator
\texttt{GEN\_SUSPENDED}, and detaches the frame from the thread state, returning control while
the frame stays alive on the heap. \texttt{.send}/\texttt{next} re-attach and resume at
\texttt{f\_lasti}. That suspendable heap frame is what \texttt{yield from} (Vol III, Ch 9) and
native coroutines (Vol III, Ch 11) generalize.

\section{Performance and anti-patterns}

\begin{itemize}
  \item \textbf{Decorator call overhead.} Each decoration adds a Python-level call per
    invocation; in hot paths a thin forwarding wrapper can dominate. Measure; consider
    \texttt{functools.cache} (a decorator that \emph{removes} work) or inlining.
  \item \textbf{Always \texttt{functools.wraps}.} Missing metadata breaks
    \texttt{inspect.signature}, \texttt{pydoc}, framework route detection, and \texttt{typing}
    tools; \texttt{\_\_wrapped\_\_} lets tools unwrap to the original.
  \item \textbf{Context managers over \texttt{try}/finally} for anything reused: composable
    (\texttt{ExitStack}), nestable, harder to get wrong. Reserve bare \texttt{try}/finally for
    one-off local cleanup.
  \item \textbf{Prime your coroutines.} Sending a non-\texttt{None} value to a just-started
    generator raises \texttt{TypeError}; new code should prefer \texttt{async}/\texttt{await}.
  \item \textbf{Don't swallow exceptions silently} in \texttt{\_\_exit\_\_}: suppress narrowly
    (check \texttt{exc\_type}) or use \texttt{contextlib.suppress} with explicit types.
\end{itemize}

\section{Summary and cross-references}

\begin{itemize}
  \item \textbf{Decorators} are compile-time sugar for \texttt{f = d(f)}, applied
    \textbf{bottom-up}; parameterized decorators are factories; class decorators (PEP 3129)
    apply the same rule to classes. Always use \textbf{\texttt{functools.wraps}}.
  \item \textbf{Context managers} are Python's \textbf{deterministic cleanup}
    primitive---RAII on a guaranteed \texttt{try}/finally; \texttt{\_\_exit\_\_} returning
    truthy suppresses the exception.
  \item \textbf{Generator coroutines} (PEP 342) add \texttt{send}/\texttt{throw}/\texttt{close},
    making a generator a \textbf{suspendable, resumable frame}---the origin of
    \texttt{async}/\texttt{await}.
  \item Modern bytecode for all three differs from pre-3.11 listings; read \texttt{dis}.
\end{itemize}

\textbf{Cross-references.} First-class functions, frames, \texttt{f\_lasti} $\rightarrow$
Chapter 1. Reference counting vs. deterministic cleanup $\rightarrow$ Chapter 2. Descriptors
$\rightarrow$ Chapter 4. \texttt{yield from} and native coroutines $\rightarrow$ Vol III, Ch 9
and Ch 11. \texttt{dataclass} class decorators $\rightarrow$ Vol IV. Full \texttt{contextlib}
$\rightarrow$ Vol X. \textbf{CPython container internals} (archived source §5.4) $\rightarrow$
Vol XII.
